\chapter{Supplementary Hardware Context}\label{app:hardware-context}

This appendix collects hardware/economic background and extended node-path
modeling details that support interpretation but are not required for reading
the main optimisation narrative.

\section{GA100 Yield Sensitivity (TSMC N7, Approximate)}
\label{sec:ga100-yield-appendix}

GA100 is a very large monolithic die (about \SI{826}{\milli\meter\squared})
with 128 physical SM partitions in the full design. Shipping A100 SXM4-40GB
parts expose 108 active SMs (20 disabled)~\cite{nvidia2020a100,nvidia-ampere-in-depth}.
This section gives a first-order yield sensitivity estimate to provide economic
context for why large dies are commonly harvested/binned.

Using the standard first-order Poisson defect model for full-die yield,
\[
  Y_{\mathrm{full}}(A,D_0)\approx e^{-A D_0},
\]
where $A$ is die area in cm$^2$ and $D_0$ is defect density in
defects/cm$^2$. For GA100, $A=8.26$.

To approximate harvesting allowances, we also model the 128 physical SM
partitions as equal critical regions with area $a_{\mathrm{SM}}=A/128$.
Under the same Poisson assumption, one partition survives with probability
\[
  q(D_0)=e^{-a_{\mathrm{SM}}D_0},
\]
so the number of broken partitions $K$ is approximated as
$K\sim\mathrm{Binomial}(128,\,1-q)$. The corresponding harvested-yield curves are
\[
  Y_{\le k}(D_0)=\Pr[K\le k]
  =\sum_{i=0}^{k}\binom{128}{i}\,(1-q)^i\,q^{128-i},
\]
for $k\in\{1,2\}$. The case $k=0$ recovers $Y_{\mathrm{full}}$.

Absolute TSMC N7 defect-density time series are not publicly released in full.
However, industry-reported anchor values from public disclosures place
$D_0\approx 0.33$ defects/cm$^2$ in a pre-MP reference regime and
$D_0\approx 0.09$ defects/cm$^2$ for N7 around three quarters after high-volume
manufacturing ramp~\cite{tomshardware2024tsmc-defect-density,anandtech2020tsmc-n5}.
These two regimes are retained as the marker points in
\cref{fig:ga100-n7-yield-curve}.

\begin{figure}[H]
  \centering
  \IfFileExists{images/ga100_n7_yield_curve.pdf}{
    \includegraphics[width=0.9\textwidth]{ga100_n7_yield_curve}
  }{
    \fbox{
      \parbox{0.86\textwidth}{
        Generate this figure with
        \texttt{python3 scripts/generate\_ga100\_n7\_yield\_curve.py}.
      }
    }
  }
  \caption{First-order defect-limited yield sensitivity for GA100 area
  (\SI{826}{\milli\meter\squared}) using three simple models: full-die
  ($0$ broken partitions), at most $1$ broken partition out of $128$, and at
  most $2$ broken partitions out of $128$. Marker points use industry-reported
  N7 anchor values ($D_0\approx 0.33$ and $D_0\approx 0.09$).}
  \label{fig:ga100-n7-yield-curve}
\end{figure}

At those two anchor points, the full-die model gives
\[
  Y_{\mathrm{full}}(D_0=0.33)\approx 6.6\%,\qquad
  Y_{\mathrm{full}}(D_0=0.09)\approx 47.5\%.
\]
Under the same simplified partition model, allowing at most one broken
partition raises the predicted values to about $24.6\%$ and $83.0\%$,
respectively; allowing at most two broken partitions gives about $49.3\%$ and
$96.1\%$.
This should not be interpreted as NVIDIA's product yield: shipping yield is
higher due to redundancy, harvesting (disabled units), and performance binning.
The model is used here only to quantify the sensitivity of large-die economics
to defect density.

\paragraph{Area increment sanity check.}
For a die-area increase $\Delta A$, Poisson yield scales as
\[
  \frac{Y(A+\Delta A,D_0)}{Y(A,D_0)}=e^{-D_0\Delta A}.
\]
With $\Delta A=0.1$\,cm$^2$, this multiplier is about $0.968$ at
$D_0=0.33$ and $0.991$ at $D_0=0.09$, i.e.\ a \SI{3.2}{\percent} and
\SI{0.9}{\percent} relative yield drop, respectively.

\section{Extended Host--Device Path Cost Model}
\label{sec:host-device-path-model}

For a transfer of size $V$ bytes, we use a first-order latency+bandwidth model
in the standard Hockney form~\cite{hockney1988parallel}:
\[
  T_{\mathrm{H2D}}(\pi, V)
  \approx
  T_0(\pi) + \frac{V}{B_{\mathrm{eff}}(\pi)},
\]
where $\pi$ is the path class (local, same-socket remote, or cross-socket
remote), $T_0$ is the fixed setup/latency term, and $B_{\mathrm{eff}}$ is the
effective end-to-end bandwidth.

To make the topology penalty explicit, we write
\[
  T_0(\pi) = T_0^{\mathrm{local}} + n_{\mathrm{IF}}(\pi)\,\tau_{\mathrm{IF}},
\]
with $n_{\mathrm{IF}}$ the number of additional Infinity Fabric segments and
$\tau_{\mathrm{IF}}$ an average per-segment latency penalty. This gives the
expected ordering
\[
  T_{\mathrm{local}} < T_{\mathrm{same\ socket\ remote}} < T_{\mathrm{cross\ socket}}.
\]

\paragraph{Concrete mapping for this node.}
Using the measured affinity in \cref{tab:gpu-numa-affinity}, the path classes
can be made explicit. For example, for GPU\,1 (local NUMA\,1):
\begin{itemize}
  \item \textbf{Local path} ($\pi=\pi_{\mathrm{local}}$):
    thread and pinned host buffer on NUMA\,1 (cores 16--31),
    modeled with $n_{\mathrm{IF}}=0$.
  \item \textbf{Same-socket remote path} ($\pi=\pi_{\mathrm{ss}}$):
    thread/buffer on NUMA\,3 (same socket, cores 48--63),
    modeled with one additional fabric segment ($n_{\mathrm{IF}}\approx 1$).
  \item \textbf{Cross-socket path} ($\pi=\pi_{\mathrm{xs}}$):
    thread/buffer on NUMA\,5 or NUMA\,7 (other socket),
    modeled with inter-socket traversal plus local fabric
    ($n_{\mathrm{IF}}\approx 2$ in this simplified model).
\end{itemize}
The same local/same-socket/cross-socket classification applies to GPU\,0, GPU\,2,
and GPU\,3 by replacing NUMA IDs with their entries from
\cref{tab:gpu-numa-affinity}. The practical expectation is
\[
  B_{\mathrm{eff}}(\pi_{\mathrm{local}})
  >
  B_{\mathrm{eff}}(\pi_{\mathrm{ss}})
  >
  B_{\mathrm{eff}}(\pi_{\mathrm{xs}}),
\]
and correspondingly higher transfer time for the same $V$.

\paragraph{Cost instantiation for this architecture.}
For topology-aware staging costs, we model a sequence of $N$ host-device
transfers with total volume $V_{\mathrm{tot}}=\sum_{i=1}^{N}V_i$ as
\[
  T_{\mathrm{stage}}(\pi)
  =
  \sum_{i=1}^{N} T_{\mathrm{H2D}}(\pi,V_i)
  \approx
  N\,T_0(\pi)+\frac{V_{\mathrm{tot}}}{B_{\mathrm{eff}}(\pi)}.
\]
With the path classes used in this node model
($n_{\mathrm{IF}}(\pi_{\mathrm{local}})=0$,
 $n_{\mathrm{IF}}(\pi_{\mathrm{ss}})=1$,
 $n_{\mathrm{IF}}(\pi_{\mathrm{xs}})=2$), this becomes
\begin{align*}
  T_{\mathrm{stage}}(\pi_{\mathrm{local}})
  &\approx
  N\,T_0^{\mathrm{local}}+\frac{V_{\mathrm{tot}}}{B_{\mathrm{eff}}(\pi_{\mathrm{local}})},\\
  T_{\mathrm{stage}}(\pi_{\mathrm{ss}})
  &\approx
  N\!\left(T_0^{\mathrm{local}}+\tau_{\mathrm{IF}}\right)
  +\frac{V_{\mathrm{tot}}}{B_{\mathrm{eff}}(\pi_{\mathrm{ss}})},\\
  T_{\mathrm{stage}}(\pi_{\mathrm{xs}})
  &\approx
  N\!\left(T_0^{\mathrm{local}}+2\tau_{\mathrm{IF}}\right)
  +\frac{V_{\mathrm{tot}}}{B_{\mathrm{eff}}(\pi_{\mathrm{xs}})}.
\end{align*}

Relative to local placement, the predicted penalties are therefore
\begin{align*}
  \Delta T_{\mathrm{ss}}
  &=
  T_{\mathrm{stage}}(\pi_{\mathrm{ss}})
  -T_{\mathrm{stage}}(\pi_{\mathrm{local}})
  \approx
  N\,\tau_{\mathrm{IF}}
  +V_{\mathrm{tot}}
   \!\left(
     \frac{1}{B_{\mathrm{eff}}(\pi_{\mathrm{ss}})}
     -\frac{1}{B_{\mathrm{eff}}(\pi_{\mathrm{local}})}
   \right),\\
  \Delta T_{\mathrm{xs}}
  &=
  T_{\mathrm{stage}}(\pi_{\mathrm{xs}})
  -T_{\mathrm{stage}}(\pi_{\mathrm{local}})
  \approx
  2N\,\tau_{\mathrm{IF}}
  +V_{\mathrm{tot}}
   \!\left(
     \frac{1}{B_{\mathrm{eff}}(\pi_{\mathrm{xs}})}
     -\frac{1}{B_{\mathrm{eff}}(\pi_{\mathrm{local}})}
   \right).
\end{align*}

\paragraph{How parameters are obtained in practice.}
For each path class $\pi$, measure transfer times for multiple sizes
$V\in[V_{\min},V_{\max}]$ and fit
\[
  T(V)=a_\pi + b_\pi V,
  \qquad
  T_0(\pi)=a_\pi,\quad B_{\mathrm{eff}}(\pi)=\frac{1}{b_\pi}.
\]
This linear fit is the standard way to instantiate latency+bandwidth models in
system-performance analysis~\cite{hockney1988parallel,jain1991art}.
Once fitted, the above expressions directly provide predicted staging overhead
for any $(N,V_{\mathrm{tot}})$ in the workload.

\paragraph{Per-GPU path classes.}
The local and same-socket ``remote'' NUMA domains used for orchestration are:
\begin{table}[H]
  \centering
  \small
  \begin{tabular}{l l l l}
    \toprule
    GPU & $\pi_{\mathrm{local}}$ & $\pi_{\mathrm{ss}}$ (used) & $\pi_{\mathrm{xs}}$ (used) \\
    \midrule
    GPU\,0 & NUMA\,3 & NUMA\,1 & NUMA\,5 or NUMA\,7 \\
    GPU\,1 & NUMA\,1 & NUMA\,3 & NUMA\,5 or NUMA\,7 \\
    GPU\,2 & NUMA\,7 & NUMA\,5 & NUMA\,1 or NUMA\,3 \\
    GPU\,3 & NUMA\,5 & NUMA\,7 & NUMA\,1 or NUMA\,3 \\
    \bottomrule
  \end{tabular}
  \caption{Path-class mapping used for topology-aware cost evaluation.}
  \label{tab:gpu-path-class-mapping}
\end{table}

For each row in \cref{tab:gpu-path-class-mapping}, the same
$n_{\mathrm{IF}}=\{0,1,2\}$ model applies, so the transfer-cost expressions
above are directly reusable across all four GPUs.
